\documentclass[english,submit]{ipsj}
% IPSJ official class. XeLaTeX + xeCJK for portable Japanese typesetting.
\usepackage{fontspec,xeCJK}
\setmainfont{Latin Modern Roman}
\setsansfont{Latin Modern Sans}
\setmonofont{Latin Modern Mono}
\setCJKmainfont{Noto Serif CJK JP}[BoldFont={Noto Serif CJK JP Bold}]
\setCJKsansfont{Noto Sans CJK JP}[BoldFont={Noto Sans CJK JP Bold}]
\setCJKmonofont{Noto Sans CJK JP}
\xeCJKsetup{PunctStyle=plain,AutoFallBack=true}
% 盤面図の持駒欄の☖☗を和文フォントで組む（Latin Modern には字形がない）
\xeCJKDeclareCharClass{CJK}{"2616,"2617}
\usepackage{amsmath,amssymb,graphicx,booktabs,array,tikz,url,float,balance}
\usetikzlibrary{arrows.meta,positioning}
\makeatletter
\def\@uketsuke{}\def\@euketsuke{}\def\SHUBETUname{}
% Japanese body leading used by the IPSJ class: 13Q type on 21H baseline.
\renewcommand{\normalsize}{\@setfontsize\normalsize{9.25pt}{14.938pt}\abovedisplayskip6pt plus2pt minus2pt\belowdisplayskip6pt plus2pt minus2pt\let\@listi\@listI}
\renewcommand{\small}{\@setfontsize\small{8.53728pt}{12.80592pt}}
\makeatother
\normalsize
\pagestyle{empty}
\newcommand{\koma}{\fontsize{12.5}{13}\selectfont}
\newcommand{\mochi}{\fontsize{8.5}{10}\selectfont}
\newcommand{\weburl}[1]{\url{https://#1}}
\newcommand{\pc}{\operatorname{pc}}
\newcommand{\V}{\mathcal V}
\renewcommand{\figurename}{図}
\renewcommand{\tablename}{表}
\renewcommand{\refname}{参考文献}
\setlength{\emergencystretch}{1em}
\begin{document}
\twocolumn[{%
\begin{center}
{\LARGE\bfseries 協力詰における全駒煙詰の発見}\\[2mm]
{\large 探索手法の段階的な設計と改良による計算機支援創作}\\[3mm]
{\large 岡 圭吾\textsuperscript{1,a)}\quad Keigo Oka\textsuperscript{1,a)}}\\[3mm]
{\large\bfseries Discovery of a 40-Piece Smoke Helpmate in Shogi}\\[1mm]
{\normalsize Computer-Aided Composition through Stepwise Design and Refinement of Search Methods}
\end{center}
\vspace{3mm}
\begin{minipage}{.94\textwidth}
\small
\textbf{概要}\quad
将棋駒一組40枚をすべて盤上に置き，113手後に3枚だけを残して詰む協力詰「狼煙」を得た．
探索は，3枚の詰上がりから局面を逆算する候補生成と，出力候補の最短解が一意であることを調べる完全探索からなる．
全幅逆算の拡大が困難になったため，探索履歴から到達駒数を予測する学習誘導ビーム探索を導入した．
さらに，得られた39枚局面の解手順が合流する構造を調べ，再始動点とビーム幅を変更することで40枚に到達した．
本稿では，逆算，教師値の改良，解構造の分析，作品選択を通じた計算機支援創作の過程を報告する．
モデルの特徴量利用も分析し，候補生成の近似性と個々の作品の検証を区別する．
得られた作品は独立の実装によっても一意性を確認した．

\vspace{2mm}
\textbf{Abstract}\quad
We obtained Noroshi, a shogi helpmate that starts with all 40 pieces of a standard set on the board and ends with three pieces after a unique shortest solution of 113 plies.
Our method separates backward candidate generation from exhaustive verification of the output candidates' shortest solutions.
When full-width backward search became impractical, we introduced a learned beam evaluation that predicts the number of pieces reachable from a position.
After reaching 39 pieces, we analyzed shared suffixes of the solutions and redesigned the restart position and the allocation of beam width, leading to 40-piece positions.
We describe the resulting computer-aided composition process, including improved training targets, structural analysis of solutions, and human selection and refinement.
We also analyze feature use in the learned model and distinguish heuristic search coverage from candidate verification.
The final composition was also checked by an independent breadth-first implementation using separately developed legal-move generation.
\end{minipage}
\vspace{5mm}
}]
\thispagestyle{empty}
% 所属とメール（jsample-gpw.tex を日本語モードで組んだときの \affiliate / \author の脚注と同じ形）
\begingroup
\makeatletter
\def\thefootnote{\@arabic{\c@footnote}\protect\hphantom{\rparen}}%
\setcounter{footnote}{0}\refstepcounter{footnote}%
\footnotetext{Google (Work done in a personal capacity)}%
\def\thefootnote{\@alph{\c@footnote}\rparen}%
\setcounter{footnote}{0}\refstepcounter{footnote}%
\footnotetext{ogiekako@gmail.com}%
\makeatother
\endgroup
\setcounter{footnote}{0}

\section{はじめに}
煙詰は，多数の駒が配置された初期状態から，手順の進行とともに駒が次々に消えていき，必要最低限の枚数で詰め上がる詰将棋の総称である．
狭義には，双玉では40枚，単玉では39枚の駒すべてを盤上に置いた初形から始まるものを指し，特にこれを全駒煙と呼ぶ．
全駒煙以外の煙詰はミニ煙と呼ばれる．
その一号局は，1755年の伊藤看寿『将棋図巧』第99番である\cite{zuko}．
2作目の黒川一郎「落花」が現れるまでの約200年間，煙詰はこの1局だけであった\cite{toybox56}．
その後は多くの作品が作られ，都煙・小駒煙・貧乏煙などを含めて2015年の時点で338作が数えられている\cite{tetsu2015}．
一方，攻方と受方が協力して詰みに向かう協力詰の煙詰は，従来実現不可能と言われてきた\cite{wfp218}．
2004年には，39枚（双玉なら40枚）から3枚にするのは「おそらく不可能」との見方も示されている\cite{mozu2004}．
本研究では，標準の将棋駒一組40枚をすべて盤上に置き，最後に3枚を残す協力詰の作品を探索した．
40枚は，この駒一組を用いる条件下で初形の盤上駒数の上限である．

得られた作品が，図\ref{fig:noroshi}の「狼煙」（のろし）である．
初形の持駒は双方ともなく，113手の最短解が一意に存在する．
詰上がりに残るのは，双方の玉と攻方のと金の3枚である．
協力詰では，作意と同じか，それより短い別の詰手順が一つでもあれば，一意な最短解という条件を満たさなくなる．
盤上の駒を増やすほど，狙った長い手順を保ちながら別解を排除することが難しくなる．

本稿の主題は，この作品の発見に至った探索設計である．
評価関数は，探索を進めながら作り直していった．
まず逆算探索で得た局面から評価を学習し，次に学習の教師値を改め，さらに探索結果の解手順を分析して起点を選び直した．
計算機が生成した結果を人間が調べ，その結果を次の探索設計に戻す過程を具体的に示す．
各構成の最大到達駒数は，発見過程の記録として報告する．

\begin{figure}[t]
\centering
\begin{tikzpicture}[x=5.2mm,y=-5.2mm]
\draw[line width=.4pt] (0,0) grid (9,9);
\draw[line width=.8pt] (0,0) rectangle (9,9);
\node[font=\footnotesize] at (0.5,-.55) {9};
\node[font=\footnotesize] at (9.5,0.5) {一};
\node[font=\footnotesize] at (1.5,-.55) {8};
\node[font=\footnotesize] at (9.5,1.5) {二};
\node[font=\footnotesize] at (2.5,-.55) {7};
\node[font=\footnotesize] at (9.5,2.5) {三};
\node[font=\footnotesize] at (3.5,-.55) {6};
\node[font=\footnotesize] at (9.5,3.5) {四};
\node[font=\footnotesize] at (4.5,-.55) {5};
\node[font=\footnotesize] at (9.5,4.5) {五};
\node[font=\footnotesize] at (5.5,-.55) {4};
\node[font=\footnotesize] at (9.5,5.5) {六};
\node[font=\footnotesize] at (6.5,-.55) {3};
\node[font=\footnotesize] at (9.5,6.5) {七};
\node[font=\footnotesize] at (7.5,-.55) {2};
\node[font=\footnotesize] at (9.5,7.5) {八};
\node[font=\footnotesize] at (8.5,-.55) {1};
\node[font=\footnotesize] at (9.5,8.5) {九};
\node[inner sep=0pt] at (8.5,0.5) {\koma と};
\node[inner sep=0pt] at (6.5,1.5) {\koma 王};
\node[inner sep=0pt,rotate=180] at (8.5,1.5) {\koma 歩};
\node[inner sep=0pt] at (4.5,2.5) {\koma 銀};
\node[inner sep=0pt] at (8.5,2.5) {\koma と};
\node[inner sep=0pt] at (3.5,3.5) {\koma と};
\node[inner sep=0pt,rotate=180] at (4.5,3.5) {\koma 歩};
\node[inner sep=0pt] at (5.5,3.5) {\koma と};
\node[inner sep=0pt] at (6.5,3.5) {\koma 歩};
\node[inner sep=0pt,rotate=180] at (7.5,3.5) {\koma と};
\node[inner sep=0pt,rotate=180] at (8.5,3.5) {\koma 桂};
\node[inner sep=0pt] at (2.5,4.5) {\koma 角};
\node[inner sep=0pt,rotate=180] at (3.5,4.5) {\koma 歩};
\node[inner sep=0pt] at (4.5,4.5) {\koma 歩};
\node[inner sep=0pt] at (5.5,4.5) {\koma 歩};
\node[inner sep=0pt,rotate=180] at (7.5,4.5) {\koma 歩};
\node[inner sep=0pt] at (3.5,5.5) {\koma 角};
\node[inner sep=0pt,rotate=180] at (4.5,5.5) {\koma 玉};
\node[inner sep=0pt,rotate=180] at (5.5,5.5) {\koma と};
\node[inner sep=0pt,rotate=180] at (6.5,5.5) {\koma と};
\node[inner sep=0pt] at (8.5,5.5) {\koma 歩};
\node[inner sep=0pt] at (1.5,6.5) {\koma 飛};
\node[inner sep=0pt] at (2.5,6.5) {\koma 飛};
\node[inner sep=0pt,rotate=180] at (3.5,6.5) {\koma 銀};
\node[inner sep=0pt,rotate=180] at (4.5,6.5) {\koma 杏};
\node[inner sep=0pt,rotate=180] at (5.5,6.5) {\koma 金};
\node[inner sep=0pt] at (6.5,6.5) {\koma 香};
\node[inner sep=0pt] at (7.5,6.5) {\koma 桂};
\node[inner sep=0pt] at (8.5,6.5) {\koma 金};
\node[inner sep=0pt] at (2.5,7.5) {\koma 銀};
\node[inner sep=0pt] at (3.5,7.5) {\koma 歩};
\node[inner sep=0pt] at (4.5,7.5) {\koma 桂};
\node[inner sep=0pt,rotate=180] at (5.5,7.5) {\koma 香};
\node[inner sep=0pt] at (6.5,7.5) {\koma 桂};
\node[inner sep=0pt] at (7.5,7.5) {\koma 歩};
\node[inner sep=0pt] at (8.5,7.5) {\koma 金};
\node[inner sep=0pt] at (0.5,8.5) {\koma 金};
\node[inner sep=0pt,rotate=180] at (1.5,8.5) {\koma と};
\node[inner sep=0pt] at (2.5,8.5) {\koma 香};
\node[inner sep=0pt] at (7.5,8.5) {\koma 銀};
\end{tikzpicture}

\caption{「狼煙」の初形（ogiekako作，『詰将棋パラダイス』2026年8月号）．協力詰113手，盤上40枚．逆向きの文字は受方の駒を表す．}
\label{fig:noroshi}
\end{figure}

\section{問題設定と関連研究}
\subsection{協力詰と本稿で扱う煙詰}
協力詰は，攻方と受方が協力して受方玉を詰めるフェアリー詰将棋である．
攻方は王手をかけ，受方は王手を解く合法手を指す．
本稿では最短の詰手順を解とし，その手順がただ一つである局面を扱う．
より長い迂回手順があっても，最短の手順が一つならこの意味で一意とする．
この規則と探索上の性質は，星らによる協力詰ソルバーの研究にも整理されている\cite{hoshi2002}．
通常の将棋と同様に，取った駒は持駒になり，再び打つことができる．
「駒が消える」とは受方の持駒になることであり，駒一組から除去されることではない．

協力詰は，1961年に『詰将棋パラダイス』67号でS・オギノ氏が「ばか詰」の名で提唱したルールである\cite{kanna}．
提唱後しばらくは顧みられなかったが，1970年の中出慶一氏の新作発表をきっかけに広まった\cite{kanna}．
直後に発表された加藤徹氏の作品によって長編の可能性も認識され，1971年には同誌に「ばか詰教室」が開設された\cite{kanna}．
1975年には加藤徹氏作の「寿限無」19447手が発表されている\cite{kanna}．

協力詰では，解がただ一つであることが重視される．
成・不成の選択や飛角香を打つ地点が一つに定まらない手を，非限定と呼ぶ．
普通詰将棋では，煙詰を含め，非限定は余詰とは区別され，許容範囲とされることが多い\cite{sumaho,komai2025}．
フェアリー詰将棋でも，以前は多少の非限定が許容されていたが，現代では唯一解を基本とすることが多く\cite{komai2025}，WFP作品展のルールも，注釈がない限り唯一解を協力詰などの完全作の条件としている\cite{wfprule}．
本稿の一意性もこの意味であり，成・不成や打つ地点だけが異なる手順も別の解として数える．
「狼煙」の最短解はこの意味で一意である（\ref{sec:verify}節）．

本稿の対象は，双方の玉を盤上に置く双玉の煙詰である．
初形では標準の駒一組（玉2，飛2，角2，金4，銀4，桂4，香4，歩18）をすべて盤上に置き，双方の持駒はなしとする．
目標は，そこからの最短の詰手順がただ一つで，その詰上がりの盤上駒数が3枚となる初形を得ることである．
駒の成りの可否，二歩，行き所のない駒，打歩詰といったルールは，すべて普通の将棋に準ずる．
手数は攻方と受方の着手を合わせて数える．

\subsection{協力詰の煙詰と計算機創作}
協力詰の煙詰については，2004年にも全駒達成の可能性が議論されている\cite{tetsu2004,mozu2004}．
25枚から3枚への到達がコメントとして報告された一方，その図は未公表であったと，後年の紹介には記されている\cite{wfp186}．
「狼煙」は，記録に残る限り，協力詰における初の全駒煙であり，『詰将棋パラダイス』2026年8月号に出題された．

著者は「狼煙」の完成に先立ち，使う駒種を限った単玉の煙詰を逆算探索で作り，第182回WFP作品展で発表した\cite{wfp218}．
歩を使わない歩なし煙（初形21枚，51手）と，金銀を使わない貧乏煙（初形31枚，75手）である．
同作品展の解説では，後者の初形31枚が単玉協力詰の煙詰として当時の最大枚数とされている\cite{wfp218}．
単玉で39枚すべてを盤上に置く全駒煙は，本稿執筆時点で未発見である．

逆算法による計算機創作には蓄積がある．広瀬正幸らは詰上がりから初形を構成する逆算法による自動創作を研究した\cite{hirose1998}．
普通詰将棋の煙詰について，広瀬稔・馬屋原は，配置駒数，駒の温存，玉の可動域などを考慮した評価関数を用いて39枚から3枚になる作品を自動生成した\cite{hirose2026}．
普通詰将棋の煙詰には『将棋図巧』以来の作品の蓄積があり，この評価関数はそこで培われた作家の知見を人手で関数にしたものである．
一方，協力詰の全駒煙には手本となる既存作品がないため，本研究では同じ役割の評価を探索の履歴から学習した．
また，黒瀧は，逆算の選択を深層学習によって誘導する詰将棋創作を報告している\cite{kurotaki2025}．
これらが評価関数や生成手法そのものを主題とするのに対し，本研究は，手本となる作品のない協力詰の全駒煙について，評価の学習，解構造の分析，作品の選択を含む創作の過程全体を扱う．

解図の側では，星らが協力詰をOR木として扱い，反復深化探索と局面表を組み合わせた手法を示した\cite{hoshi2002}．
普通詰将棋の攻方と受方の対立に対応するAND/OR木とは異なり，協力詰では，双方の合法手すべてを，詰みへ向かう可能性として調べる．
本研究で用いた自作プログラムfmrsは，逆算候補の絞り込みにはキャッシュ付き深さ制限DFSを，最終局面の解図には幅優先探索によるソルバーを用いる\cite{fmrs}．

\section{逆算による候補生成と一意性の検証}
\subsection{二つの探索を分離する}
探索システムは，詰上がりから初形の方向に進む逆算探索と，候補局面から詰みに向かって解を調べる検証探索からなる．
逆算探索は作品の候補を作る役割を持ち，計算量を抑えるために局面を間引いてよい．
検証探索は候補が作品として成立するかを判定する役割を持ち，常に全合法手を調べる．

逆算では，合法な着手を逆向きに戻して一つ前の局面を生成する．
移動を戻すほか，駒取りを戻して盤上に駒を復元する場合や，成りを戻す場合がある．
持駒を打った手を戻すと，盤上の駒は持駒に戻る．
このため盤上駒数は，逆算の1手ごとに増える場合も減る場合もある．
逆算は詰みからの手数ごとに層をなして進む．
各層では，生成した局面の重複を除き，駒数や駒の種類に関する探索用の制約（\ref{sec:constraints}節）を適用したうえで，一意性の判定を通った局面を次の層の候補とする．
この各層の候補の集合をフロンティアと呼ぶ．
フロンティアの局面をすべて次の層へ進める探索を全幅探索，評価値などに基づいて決まった数（ビーム幅）の局面だけを残して次の層へ進める探索をビーム探索と呼ぶ．

候補局面について，逆算で付け加えた既知の手順の長さを $D$ とする．
検証では，最短の詰手順の長さがちょうど $D$ で，その手順が1通りであることを調べる．
解が2通り見つかればその候補は棄却できるため，解数は0，1，2以上を区別すればよい．
解は手順として数えるので，同じ局面に異なる手順で至る場合はそれぞれを別の解とする．
fmrsでは，残り深さに関する情報を持つ局面表と深さ制限DFSを用いて，この判定を繰り返す．
局面表は判定結果の再利用にだけ使うので，局面表を捨てて探索を再開しても判定結果は同じで，計算量だけが変わる．

駒構成の制約とビームによる選択は逆算探索にだけ適用し，検証探索は常に全合法手を対象とする．
こうすることで，逆算で捨てた枝にある余詰も検証で検出できる．
また，候補を絞る途中では判定を速くするため，攻方の金と，金と同じ動きをする成香・成桂・成銀を，受方の持駒の歩と入れ替えてと金に置き換える正規化を用いた（歩が足りなければ香・桂・銀・金の順に使う）．
同じ働きの駒の違いだけで別々に判定していた局面をまとめ，判定結果の再利用を増やすためである．
この置き換えで解が変わる場合が稀にあるので，最終候補は駒種と持駒を元のまま扱うソルバーで改めて検証する．

\subsection{詰上がりと探索起点の選択}
全駒煙の探索に用いた詰上がりは，1一玉，3二王，2二とからなる双玉3枚の形である（図\ref{fig:mate}）．
受方の残りの駒は持駒にあり，攻方の持駒はない．
1一玉は2二との王手を受け，1二，2一への移動はと金に，2二のと金を取る手は3二王に阻まれる．

この詰上がりは，探索の結果を見て人手で選んだものである．
単玉の少数駒煙の探索を経て，双玉なら収束の形を増やせると考え，双玉の3枚詰上がりの候補すべてから逆算を行った．
得られたログで逆算の広がりと到達駒数を目視し，有望と判断した収束を選んで計算を集中した．
途中では，有望と判断した中間局面を人手で起点として与えた探索も行っている．

\begin{figure}[t]
\centering
% generated by sfen2tikz.py -- do not edit
% sfen: 8k/6K+P1/9/9/9/9/9/9/9 w 2r2b4g4s4n4l17p 114
\begin{tikzpicture}[x=-5.2mm,y=-5.6mm]
\foreach \i in {1,...,8} {
  \draw[line width=0.1mm] (\i,0) -- (\i,9);
  \draw[line width=0.1mm] (0,\i) -- (9,\i);}
\draw[line width=0.4mm] (0,0) rectangle (9,9);
\foreach \p in {(3,3),(6,3),(3,6),(6,6)} \fill \p circle [radius=0.3mm];
\node[inner sep=0pt,rotate=180] at (0.5,0.5) {\koma 玉};
\node[inner sep=0pt] at (2.5,1.5) {\koma 王};
\node[inner sep=0pt] at (1.5,1.5) {\koma と};
\node[inner sep=0pt] at (9.75,0.4) {\mochi ☖};
\node[inner sep=0pt] at (9.75,1.22) {\mochi 残};
\node[inner sep=0pt] at (9.75,2.04) {\mochi り};
\node[inner sep=0pt] at (9.75,2.86) {\mochi 全};
\node[inner sep=0pt] at (9.75,3.6799999999999997) {\mochi 部};
\node[inner sep=0pt] at (-0.75,0.4) {\mochi ☗};
\node[inner sep=0pt] at (-0.75,1.22) {\mochi な};
\node[inner sep=0pt] at (-0.75,2.04) {\mochi し};
\end{tikzpicture}

\caption{「狼煙」の詰上がり．受方1一玉，攻方3二王と2二とが残る．盤上は3枚だが，取られた駒は受方の持駒となっている．}
\label{fig:mate}
\end{figure}

\subsection{全幅逆算の限界と駒構成の制約}
\label{sec:constraints}
逆算のフロンティアは深さとともに大きくなる．
初期の全幅探索は，歩の割合60\%以上・7筋以内という最終実行より強い制約の下で，768\,GBメモリの計算機を用いて詰みから49手まで行い，最大22枚を得た．
フロンティアは37手で約406万局面，49手で約2975万局面であり，128\,GBの計算機では38〜39手付近でメモリが不足した．
同じ記録で，37手で最多の18枚となった352局面の解経路をたどると，11〜37手の各手数でこれらの経路上にある局面はフロンティアの0.002〜0.1\%で，途中では1〜2局面に合流していた（金類の正規化後の局面数）．
詰みから11手のフロンティアのうち，より深い手数で最多駒数となる局面につながるものは2.6\%（31手では0.02\%）であった．
最多駒数につながる局面はフロンティアのごく一部であり，それを事前に見分けられれば探索を大きく絞れることになる．

一意性の検証については，詰みから31手までの計測で，処理の約97.5\%が一意な局面について全合法手を調べ尽くす部分であった．
別解のある局面は平均1.15手の展開で棄却できており，手の順序の改善で短縮できる余地は小さかった．
そこで40枚まで延ばすため，検証を速くする代わりに，全幅探索からビーム探索へ移り，展開する局面を選ぶことにした．

40枚を得た探索では，標準の駒一組を超えないことに加え，歩の割合と飛角・香桂の使用時期を制限した．
歩の割合は，詰みから6手以上の範囲で，盤上と攻方持駒にある玉以外の駒を分母として44\%以上とする．
成駒は元の駒種に含めて数える．
飛角系は使用駒数31枚から，香桂系は8枚から許し，以後，使用駒数が2枚増えるごとに許容枚数を1枚増やす．
ここで使用駒数は，盤上と攻方持駒の合計をいう．
また，受方の持駒に歩が残っている間は，攻方の自陣（7〜9段）に成香・成桂・成銀を置かないこととした．
同じ働きのと金で代わりが利く場合に，と金を優先するためである．
これらは，作品らしい駒種の残し方へ探索を誘導するために設けた探索上の制約である．

\section{探索履歴から学ぶビーム評価}
\subsection{最大到達駒数を教師値とする}
全幅探索から，各層で無作為に選んだ $10^5$ 局面だけを残すビーム探索へ移ると，初形30枚の煙詰まで到達した．
さらに，有望な局面を多く残すため，ある局面から逆算を続けたときに到達できる最大駒数を予測する評価関数を学習した．
目的は40枚に届く系列をビームに残すことなので，評価もその見込みだけを表すように設計した．

局面 $s$ の盤上の駒の枚数（玉を含む）を $\pc(s)$（piece count）と書く．
また，$s$ から逆算をさらに進めて得られる前局面のうち，探索で実際に生成し記録したものの集合を $R(s)$ とする．
$s$ 自身も含め，教師値を
\begin{equation}
 y(s)=\max_{t\in R(s)\cup\{s\}}\pc(t)
 \label{eq:target}
\end{equation}
と定める．
これは，探索で実際に見つかった範囲で，その局面からどこまで駒数を増やせたかを表す値である．

当初は，得られた解手順をたどり，その手順上で見た駒数を各局面の教師値としていた．
ところが，ある局面に複数の逆算先があると，たどらなかった枝で発見済みの大きな駒数が教師値に反映されない．
そこで，生成した局面の親子関係を記録し，探索済みグラフの全前局面から式\eqref{eq:target}を集計するよう変更した．
深い側から値を伝播すれば，直接の逆算先を $C(s)$ として
\[
 y(s)=\max\bigl(\pc(s),\ \max_{t\in C(s)}y(t)\bigr)
\]
で計算できる．
この変更は，すでに探索で得ていた情報を教師値に漏れなく反映させるためのものである．
変更前後のモデルを，同じ残り手数・同じ盤上駒数の局面どうしの順位相関（Spearman）で比べると，0.264から0.494に上がった（762万局面と1570万の親子関係から作った学習資料で，解経路単位に分けた交差検証による）．
変更前の教師値では，目的関数を順位学習に変えても同じ指標は0.32程度にとどまっており，予測の精度を決めていたのは教師値の質だったと考えられる．
なお，これは学習時の予測精度の指標である．

\subsection{特徴量と選択の多様性}
局面から表\ref{tab:features}の85次元の特徴量を抽出し，勾配ブースティング木による回帰を用いた．
玉の周囲の空きマス，利きの密度，駒の広がりなどは，逆算の余地を数値として与えるための特徴である．
攻方の王手の数を数える1特徴は計算が重いため，通常の設定では0に固定した．
公開された最終モデルは300本の木からなる\cite{fmrs}．

\begin{table}[t]
\centering\small
\caption{特徴量の構成．王手数の1特徴は通常設定で無効．}
\label{tab:features}
\begin{tabular}{lrp{37mm}}
\toprule
分類 & 次元 & 例 \\
\midrule
進行度・駒総数 & 6 & 残り手数，盤上・持駒総数 \\
盤上の駒種 & 28 & 14駒種の先後別枚数 \\
持駒の駒種 & 14 & 7駒種の先後別枚数 \\
玉の位置・包囲 & 10 & 玉座標，近傍の駒と利き \\
盤面の広がり & 8 & 使用範囲，歩の筋数 \\
玉の自由度 & 8 & 空きマス，2手先の余地 \\
配置の散らばり & 10 & 重心距離，外接矩形 \\
攻方の王手数 & 1 & 合法な王手の数 \\
\bottomrule
\end{tabular}
\end{table}

評価値の上位だけを残すと，よく似た局面にビームが偏ることがあった．
そこで，予測値に乱択を加え，駒数ごとに選択枠を分ける方式も試した．
具体的には，局面 $s$ の予測値 $f(s)$ に対して独立なGumbel乱数 $g_s$ を用い，
\[
 f(s)+Tg_s
\]
を選択のスコアとする．温度 $T$ が0なら予測値による貪欲な選択になり，$T$ を大きくすると乱択の比重が増す．
さらに，試行した選別方式の一つでは，駒数の異なる局面を順番に採用し，一部の駒数帯が候補枠を占有するのを抑える．
全幅探索の記録と同じ起点・制約で，教師値改良前のモデルを幅$5\times10^4$のビームに用いて比べると，評価値の上位だけを残す選択は21枚で行き詰まった．
無作為な選択は71手で28枚に達し，$T=5$の乱択を加えた選択は，49手で全幅探索と同じ22枚に達し，67手で28枚に達した．
評価の精度だけでなく，異なる逆算経路を残すことが必要だった．

\subsection{モデルが利用した特徴の分析}
どの特徴が使われたかを把握するため，公開モデルの木構造を調べた．
エクスポートされた木には学習標本数が残っていないため，木の形だけから計算できる次の量を用いた．
分岐ノード $v$ の左右の部分木について，葉の予測値の単純平均を $\bar a_L,\bar a_R$，ノード以下の葉数を $m_v$，その木の全葉数を $m_T$ とし，
\[
 q(v)=|\bar a_L-\bar a_R|\,m_v/m_T
\]
をその分岐の特徴に加算した．全特徴の合計で正規化した値を，構造的寄与率と呼ぶ．

\begin{table}[t]
\centering\small
\caption{最終モデルの構造的寄与率の上位．モデルが各特徴を分岐にどれだけ使ったかを表す．}
\label{tab:importance}
\begin{tabular}{lr}
\toprule
特徴 & 寄与率（\%）\\
\midrule
盤上駒数 & 22.71\\
詰みまでの残り手数 & 7.80\\
玉と駒重心のチェビシェフ距離 & 3.69\\
玉の2手先までの逃走余地 & 3.68\\
攻方の利きがある玉の近傍マス数 & 3.50\\
攻方の利きの盤面総和 & 3.18\\
玉近傍の空きマスへの平均被利き数 & 3.16\\
玉から8方向への空きマス総数 & 3.15\\
\bottomrule
\end{tabular}
\end{table}

結果を表\ref{tab:importance}に示す．盤上駒数が最大となるのは，教師値が現在の駒数以上であることとも整合する．
同じ残り手数・同じ駒数の候補を比較するときには，この二つは一定であり，玉の自由度や配置の広がりを表す特徴が候補間の違いを与える．
構造的寄与率は，学習済みモデルが各特徴をどの程度使ったかを表す．
どの特徴が40枚の発見に不可欠だったかを知るには，特徴を除いて再学習し，再探索して比べる必要がある．

\section{39枚から40枚への探索の再設計}
\subsection{到達駒数の推移}
探索構成を変更した際の到達点を表\ref{tab:milestones}にまとめる．
学習評価に乱択（$T=15$）を加えた幅$10^5$の構成で35枚，ビーム幅と選別規則を調整した構成で36枚となった．
$T=15$では，同じ設定でも到達駒数が実行ごとに3〜6枚程度ばらついた．
一方，学習導入前には，人手で選んだ中間局面から間引きせずに探索して38枚を得ていた．
探索済みグラフから教師値を集計し直したモデルを用いると，39枚に到達した．

\begin{table}[t]
\centering\small
\caption{発見過程の到達点．各行は，起点，計算量，探索制約がそれぞれ異なる実行の結果である．全幅逆算探索の行は歩の割合60\%以上・7筋以内の制約による．38枚の行は学習導入前の別実行である．}
\label{tab:milestones}
\begin{tabular}{p{59mm}r}
\toprule
探索構成 & 駒数\\
\midrule
全幅逆算探索（詰みから49手まで） &22\\
無作為ビーム（幅$10^5$） &30\\
学習評価と乱択（幅$10^5$） &35\\
幅・選別規則を調整 &36\\
人手で選んだ中間局面からの全幅探索 &38\\
全子孫から教師値を再集計した学習評価 &39\\
共通手順の分析に基づく再始動 &40\\
\bottomrule
\end{tabular}
\end{table}

表の行の間では，計算資源，中間局面の選択，探索制約，評価関数とその学習資料が同時に変わっている．
この表が示すのは，各構成で実際にその駒数の一意解を持つ局面を得たことと，その結果を次の設計に使った順序である．
個々の工夫の寄与を分けて示すには，条件をそろえた比較が要る．

\subsection{多数の初形が同じ手順へ合流する}
39枚の局面が438局得られた時点で，解手順を比較した．
すべて107手詰であり，金類の正規化を行って解手順を比較すると，詰みに至る最後の82手を共有していた．
差があるのは初形側の25手である．
438局の全手順が通る最も初形寄りの共通局面は，詰みまで82手，盤上27枚であった．
このような共通局面を合流点と呼ぶ．

解を指す方向には多くの初形が一つの手順へ合流するが，逆算の方向には一つの局面から多くの初形が分岐する．
従って，この27枚局面を新たな起点にすれば，すでに共有されている82手の部分を逆算し直さず，初形側の残りの部分に計算を集中できる．
この部分を，同じ駒構成制約と金類の正規化の下で，ビームで間引かずに探索すると，最大は39枚だった．
別に得られていた二つの39枚系列についても，それぞれの合流点から初形側を全幅で調べると，最大は39枚だった．

これは，調べた合流点と，用いた制約・正規化の設定の下での観測である．
そこで40枚の探索では，既知の39枚の系列を延ばす代わりに，より詰みに近い側から別の分岐を探す方針を採った．

\begin{figure}[t]
\centering
\begin{tikzpicture}[x=1cm,y=1cm,>=Stealth,font=\small]
\node[draw,align=center,minimum width=27mm] (mate) at (0,0) {3枚の詰上がり\\残り0手};
\node[draw,align=center,minimum width=27mm] (gate) at (0,1.3) {共通の再始動点\\14枚，残り42手};
\node[draw,align=center] (a) at (-1.75,2.85) {系列Aの合流点\\27枚，残り82手};
\node[draw,align=center] (b) at (1.75,2.85) {別の39枚系列\\B，Cの合流点};
\node[draw,align=center] (aa) at (-1.75,4.2) {438局の39枚初形\\各107手};
\node[draw,align=center] (bb) at (1.75,4.2) {B，Cの39枚初形};
\draw[->] (gate)--(mate);
\draw[->] (a)--(gate);\draw[->] (b)--(gate);
\draw[->] (aa)--(a);\draw[->] (bb)--(b);
\draw[->,dashed,bend right=38] (gate.east) to node[right,align=left,font=\footnotesize]{ここから再探索\\別の分岐を残す} (3.2,1.9);
\node[align=left,font=\footnotesize] at (-2.55,.1) {実線の矢印は\\詰みに向かう方向};
\end{tikzpicture}
\caption{金類の正規化後の39枚解の合流構造と再始動の考え方．模式図であり，辺の長さは手数を表さない．}
\label{fig:confluence}
\end{figure}

\subsection{詰みに近い共通局面から再始動する}
正規化して比較した三つの39枚系列が共通に通る最も初形寄りの局面は，詰みまで42手，盤上14枚であった．
そこでこの局面を起点に戻り，初形側でビーム幅を増やす探索を行った．
最終到達実行では96 vCPU，768\,GBメモリのクラウド計算機を用い（使用メモリのピークは68\,GB，局面表は約9億件），基本幅を100万，残り75手の位置で300万とする設定を用いた．
指定した点の間は幅を幾何的に変化させる設定とし，候補が増える深さに多くの枠を配分した．

この探索により，残り113手の位置で40枚・持駒なしの局面が2618個得られた．
探索は40枚の発見後も続け，ビームに残した候補から先へ進める局面がなくなった時点で終えた．
2618個は，このビーム実行が保持した候補から得られた40枚局面の集合である．

\section{得られた作品と検証}
\subsection{近縁な候補から作品を選ぶ}
2618個の候補は終盤の手順を広く共有しており，受方玉の軌跡，攻方が取る駒種の列（成駒を区別），受方が駒を取る位置の列を組み合わせて分類すると169類にまとまる．
成駒の種類を置き換えたものや，初形側の配置だけが異なる近縁な変化も多い．

作品の選択では，成銀・成桂・成香などの種類と枚数，成駒の位置，受方玉と攻方の駒のまとまりを調べ，候補を3局に絞った．
その後，人手で角の位置を調整し，同じように働く駒をより自然な駒種や成否へ変更して初形を整えた．
修正のたびに最短解の一意性を再検証した．
図\ref{fig:noroshi}はこの仕上げ後の発表図であり，2618候補のどれとも一致しない．

なお，初形の3七香は，取り除いて受方の持駒に移しても最短解が同じになる飾り駒である．
全駒煙の条件は40枚すべてが盤上にあることであり，飾り駒があっても条件は満たす．
一方，飾り駒のない全駒煙も得られている．
同じ基準でfmrsのソルバーを用いて各駒を一つずつ取り除いて調べると，2618候補のうち2576局は飾り駒を1枚（1五香または3七香）含んでいたが，42局には飾り駒がなかった．
\ref{sec:later}節で述べる115手の67局にも，飾り駒はなかった．
これらのうち2局（2618候補と115手から1局ずつ）は，独立の幅優先探索でも，どの駒を取り除いても最短解が変わることを確かめた．

手順の概略を見ると，序盤では銀・飛・角などを順次処理し，中盤では金や歩を使って玉を下辺から右辺へ導く．
終盤はと金と歩を消去し，図\ref{fig:mate}の3枚に至る．
本作の全手順を表\ref{tab:moves}に示す．
各段階で減る盤上駒数だけでなく，持駒を打って再利用する点も煙詰の構成に含まれている．

\begin{table*}[t]
\centering\small
\caption{「狼煙」の113手の解手順．各行は左から右に指す．数字は開始手数．「打」は駒打ち，「成」は成りを表す．}
\label{tab:moves}
\begin{tabular}{r@{\quad}l}
\toprule
1 & ６七銀\quad 同成香\quad ５七銀打\quad 同成香\quad 同飛\quad 同金\quad 同飛\quad 同と\quad ６七金打\quad ４七玉\\
11 & ５七金\quad ３七玉\quad ４七金\quad 同と\quad ４八角\quad 同と\quad 同角\quad 同玉\quad ４九香打\quad ５八玉\\
21 & ５九香打\quad ６八玉\quad ６九歩打\quad ６七玉\quad ６八歩\quad ６六玉\quad ６五と\quad 同玉\quad ６四銀成\quad 同玉\\
31 & ５四と\quad ６五玉\quad ６四と\quad ６六玉\quad ６七歩\quad 同玉\quad ６八歩打\quad 同玉\quad ６九歩打\quad ７九玉\\
41 & ８九金\quad ６九玉\quad ７九金\quad ５九玉\quad ６九金\quad ４九玉\quad ５九金\quad ３九玉\quad ４九金\quad ２九玉\\
51 & １九金\quad 同玉\quad １八金\quad ２九玉\quad ３九金\quad １八玉\quad ２九金\quad １七玉\quad １八金\quad １六玉\\
61 & １七金\quad 同玉\quad １八歩打\quad ２八玉\quad ２九歩打\quad ２七玉\quad ２八歩\quad １六玉\quad １七歩\quad １五玉\\
71 & １四と\quad 同玉\quad ２六桂打\quad 同歩\quad 同桂\quad １五玉\quad １六歩\quad ２六玉\quad ２七歩\quad ２五玉\\
81 & ２六歩\quad １四玉\quad １五歩\quad １三玉\quad １四歩\quad 同と\quad １二と\quad ２四玉\quad ２五歩\quad ３四玉\\
91 & ３五歩打\quad ４五玉\quad ４六歩打\quad ４四玉\quad ４五歩\quad ５五玉\quad ５四と\quad ４五玉\quad ４四と\quad ３五玉\\
101 & ３四と\quad ２五玉\quad ２四と\quad １五玉\quad １四と\quad 同玉\quad １三と\quad 同玉\quad １四歩打\quad １二玉\\
111 & １三歩成\quad １一玉\quad ２二と\\
\bottomrule
\end{tabular}

\end{table*}

\subsection{検証が保証する範囲}
\label{sec:verify}
作品化の過程では，神無次郎氏による独立の検討ソフトfmも用いて一意性を確認した\cite{fm}．
最終図については，本稿の再確認としてpython-shogi 1.1.1の合法手生成\cite{pythonshogi}を用いた独立の幅優先探索を実装し，最短距離と最短解数を求めた．
この探索も1解113手を返し，手順はfmrsの結果と一致した．
解図ソルバーが保証するのは解の一意性であり，初形の美的な選択は作者の判断による．

発表図のSFEN，探索コード，学習済みモデルおよび特徴分析コードは公開リポジトリで参照できる\cite{fmrs}．
本稿の図と着手表は同一のSFENと着手列から生成している．
保存された2618候補も原局面のまま再検証し，すべて1解113手，標準の駒一組による40枚初形と3枚の詰上がりを確認した．
標準駒一組の内訳と最終盤面は，着手を順に適用しても確認できる．
独立な幅優先探索では盤面・手番・双方持駒のSFENをキーとして状態を区別し，最短経路数は2以上を2に丸めて数えた．
得られた113手はすべて合法で，各攻方着手が王手であること，双方の玉が残ること，終端が盤上3枚の詰みであることも着手の逐次適用で確認した．

\subsection{共通構造と後続の系列}
\label{sec:later}
この探索で得られた40枚候補には，共通する終盤と，初形側に集中する変化という特徴がある．
探索がどれだけ異なる作品構造を見つけたかを見るため，初形の個数に加えて，手順に基づく分類を併記した．
保存された2618候補を改めて集計すると，受方玉の位置の列は全候補で一致し，攻方が取る駒種の列は146通り，受方が駒を取る位置の列は6通りであった．
これらを組み合わせた組が169類である．
この分類は各着手の時点を保持し，駒取りでない手はその旨を記録する．
さらに，手数カウンタを除く盤面・手番・双方持駒をそのまま比較しても，最後の43手は全候補で共通していた．
従って，169類は，同じ玉の追い方の上で初形側の駒の取り方が異なる変化である．
39枚段階で合流構造を調べたことも，多数の局面を得ながら実際には限られた系列を延ばしていたと気づくために役立った．

後に，起点を詰みから34手・盤上11枚の局面まで戻した探索（幅を200万から始め，残り75手で1000万，上限3000万）では，40枚局面が109手で92局，113手で6765局，115手で67局得られた（109手と113手は代表局のみ保存した）．
同じ起点から幅を小さくした別の実行（幅20万から300万）でも113手の40枚局面が1084局得られ，そのうち2618候補と重なったのは36局であった．
これらから，40枚局面は2618個のほかにも多数あり，少なくとも115手のものがあることが分かる．
ただし，これらの別系列もすべて同じ詰上がりから生成されたものである．
異なる詰上がりを含む作品群の比較は今後の課題である．

\subsection{豆腐図式の煙詰}
同じ逆算探索と検証で，盤上の駒が玉と歩・と金だけの豆腐図式の煙詰も調べた．
詰上がりを3枚（受方玉を5五で詰める双玉の都詰では4枚）とし，ビームによる間引きをせずに逆算候補が尽きるまで探索した．
攻方の王手は歩・と金による隣接王手に限られ，合駒が生じないので，解手順の全局面で盤上を玉・歩・と金に限っても，すべての候補を調べられる．
初形の最大枚数は，単玉で11枚（29手），双玉で16枚（53手），双玉の都詰で12枚（27手）であった．
初形に攻方の持駒を許して盤上と持駒の合計で数えても，最大枚数は変わらなかった．
最大局面の最短解が一意であることは，全局面についてfmrsのソルバーで，一部の局面について独立の幅優先探索で確かめた．
これらが最大であることは，逆算探索がすべての候補を尽くしたことに基づく．

\section{考察}
\subsection{人間と計算機の役割}
本事例では，計算機が合法な逆算候補を大量に生成し，その解を検証した．
人間は，詰上がりと中間起点を選び，探索の停滞を調べ，評価の教師値と探索範囲を変更した．
最後には作品の初形を選び，駒の配置と成否を整えた．
このため，本成果は，探索結果を設計に戻す計算機支援創作として位置付ける．

重要だったのは，最大駒数に加えて，得られた解の集合の構造を調べたことである．
39枚の初形が多数あっても，後半の手順が同じなら，探索していたのは一つの系列だけである．
合流点の先を全幅で調べることで，既知の終盤を維持したまま延長する方針の限界を確認し，それより詰みに近い局面から分岐を探す方針に変更できた．
これは，各局面の評価だけでは見えにくい，解の集合の性質を使った設計変更である．

\balance
\subsection{一般化に向けた課題}
本手法を別の創作課題に適用するには，少なくとも三つを検討する必要がある．
第一に，教師値である最大到達駒数は全駒煙の目標に沿うが，長手数，手筋の豊かさ，初形の自然さなどを重視する作品では目的に合わせた評価が必要になる．
第二に，手順の合流を利用するには，共通部分を固定する利点と，新しい系列を捨てる危険の両方を考える必要がある．
第三に，特定の起点と探索履歴で学習したモデルが別の収束でも役立つかは，確かめる必要がある．

評価や選別法の効果をさらに定量化するには，起点，制約，乱数種，計算予算をそろえた比較が必要である．

\section{おわりに}
協力詰において，標準の将棋駒一組40枚をすべて盤上に置き，113手後に3枚を残して詰む作品「狼煙」を得た．
最短解の一意性を検証する探索を維持した上で，候補生成を学習誘導ビーム探索に置き換え，教師値と再始動点を改良した．
39枚解が共有する手順を調べたことが，既知の系列から離れた分岐へ探索を向ける契機になった．
全駒煙の発見は，候補生成，検証，学習，解構造の分析と人間による作品化を繰り返した結果である．

\section*{謝辞}
『詰将棋パラダイス』での特別出題を担当され，進捗報告のたびに励ましをくださった神無七郎氏に感謝する．

\begin{thebibliography}{99}
\raggedright
\bibitem{zuko} 伊藤看寿：将棋図巧，第99番（1755）．
\bibitem{toybox56} TETSU：記録展示室 No.56，詰将棋おもちゃ箱．\weburl{www.ne.jp/asahi/tetsu/toybox/challenge/ct056.htm}（参照2026-10-12）．
\bibitem{tetsu2015} TETSU：看寿賞作品の煙詰率，詰将棋メモ（2015）．\weburl{toybox.tea-nifty.com/memo/2015/07/post-9f8e.html}．
\bibitem{tetsu2004} TETSU：協力詰の煙詰，詰将棋メモ（2004）．\weburl{toybox.tea-nifty.com/memo/2004/10/post_23.html}．
\bibitem{mozu2004} mozuyama：ばか煙，勝手に将棋トピックス（2004）．\weburl{mozuyama.hatenadiary.org/entry/20041024/P20041024BAKAKEMURI}．
\bibitem{wfp186} たくぼん：緩募！ナンマイダ～コンテスト，Web Fairy Paradise，No.186，p.78（2023）．\weburl{tsume-springs.com/web-fairy-paradise/WFP186.pdf}．
\bibitem{wfp218} 神無七郎：第182回WFP作品展結果，Web Fairy Paradise，No.218，pp.28--34（2026）．\weburl{tsume-springs.com/web-fairy-paradise/WFP218.pdf}．
\bibitem{kanna} 神無七郎：神無一族の氾濫，第19回（2003）．\weburl{k7ro.sakura.ne.jp/overflow/hr19.htm}（参照2026-10-12）．
\bibitem{sumaho} スマホ詰パラ：詰将棋のルール．\weburl{aws.tumepara.com/tsp/webs/rule}（参照2026-10-12）．
\bibitem{komai2025} 駒井めい：フェアリー詰将棋の細かなあれこれ，note（2025）．\weburl{note.com/meikomai_tsume/n/n33ff86cfbe6b}（参照2026-10-12）．
\bibitem{wfprule} 神無七郎：WFP作品展登場ルール一覧，Web Fairy Paradise（2026）．\weburl{tsume-springs.com/web-fairy-paradise/wfprule.pdf}（参照2026-10-12）．
\bibitem{hirose1998} 広瀬正幸，伊藤琢巳，松原仁：逆算法による詰め将棋の自動創作，人工知能学会誌，Vol.13，No.3，pp.452--460（1998）．
\bibitem{hoshi2002} 星由雄，野下浩平，柳井啓司：反復深化探索に基く協力詰将棋の解法，情報処理学会論文誌，Vol.43，No.1，pp.11--19（2002）．
\bibitem{hirose2026} 広瀬稔，馬屋原剛：独自の評価関数に基づいた煙詰の自動生成，情報処理学会研究報告，Vol.2026-GI-57，No.4，pp.1--8（2026）．
\bibitem{kurotaki2025} 黒瀧大地：深層学習を用いた逆算法による詰将棋の自動創作，第66回プログラミング・シンポジウム予稿集，pp.103--111（2025）．\weburl{prosym.org/66/program.html}．
\bibitem{fmrs} Oka, K.: fmrs，公開コード・作品・解析資料．\weburl{github.com/ogiekako/fmrs}（参照2026-10-09）．資料の参照版はmainの546167f，求解に用いた公開版はgh-pagesの24a372d．作品再生：\weburl{ogiekako.github.io/fmrs/noroshi}．
\bibitem{pythonshogi} gunyarakun: python-shogi, version 1.1.1．\weburl{github.com/gunyarakun/python-shogi}．
\bibitem{fm} 神無次郎：fm．\weburl{k7ro.sakura.ne.jp/fm.html}（参照2026-10-09）．
\end{thebibliography}
\end{document}
